\documentclass[a4paper,11pt]{article}

\usepackage[utf8]{inputenc}
\usepackage[T1]{fontenc}
\usepackage[brazil]{babel}
\usepackage[margin=2.5cm]{geometry}
\usepackage{booktabs}
\usepackage{tabularx}
\usepackage{enumitem}
\usepackage{parskip}
\usepackage{tikz}
\usetikzlibrary{shapes.geometric, arrows.meta, positioning, fit}
\usepackage[colorlinks=true, linkcolor=black, urlcolor=blue]{hyperref}

\title{\textbf{Fundação Toque --- Sistema Web de Gestão}\\[6pt]
\large Estrutura, Arquitetura e Banco de Dados}
\author{Equipe RootTech --- IFSP Câmpus Araraquara \\ \small Disciplina de Projeto de Extensão}
\date{\today}

\begin{document}

\maketitle

\section{Estrutura do Projeto}

O repositório do projeto está organizado da seguinte forma:

\begin{verbatim}
fundacao-toque/
|-- src/                        Aplicação web (React + TypeScript)
|   |-- features/               Telas organizadas por funcionalidade
|   |   |-- auth/               Login e controle de sessão (AuthContext)
|   |   |-- alunos/             Listar, cadastrar e detalhar alunos
|   |   |-- frequencia/         Chamada e contatos por faltas
|   |   `-- usuarios/           Gestão de perfis (Direção)
|   |-- services/               Regras de negócio e acesso ao Firestore
|   |   |-- studentsService.ts  Operações de alunos
|   |   `-- attendanceService.ts Operações de frequência
|   |-- lib/                    Integrações e utilitários
|   |   |-- firebase.ts         Configuração do SDK do Firebase
|   |   `-- masks.ts            Máscaras de validação (CPF, telefone etc.)
|   |-- types/                  Definições de tipos TypeScript
|   `-- App.tsx                 Roteamento e telas principais
|-- prototipo/                  Protótipo estático (usado como base)
|   |-- alunos/                 Listar, cadastrar e detalhar alunos
|   |-- frequencia.html         Chamada com presente/ausente
|   |-- evolucao.html           Registro de evolução pedagógica
|   |-- atividades.html         Registro diário + upload de fotos
|   |-- relatorios.html         Prestação de contas (PDF)
|   |-- usuarios.html           Gestão de perfis (Direção)
|   `-- Documento_de_Requisitos_Fundacao_Toque.md
|-- scripts/
|   `-- seed-firestore.mjs      População/limpeza do banco (dados fictícios)
|-- firestore.rules             Regras de segurança (controle de acesso)
|-- firestore.indexes.json      Índices compostos declarados
|-- firebase.json               Configuração do Firebase CLI
|-- modelagem.md                Documentação da modelagem do banco
`-- package.json                Dependências e scripts
\end{verbatim}

\subsection*{Telas do sistema (perfis de acesso)}

\begin{center}
\begin{tabularx}{\textwidth}{lX}
\toprule
\textbf{Perfil} & \textbf{Telas e funcionalidades} \\
\midrule
Direção & Acesso completo; gestão de usuários; relatórios de prestação de contas \\
Pedagogia & Cadastro e acompanhamento pedagógico, evolução, oficinas \\
Administrativo & Cadastro de alunos, frequência, documentação \\
\bottomrule
\end{tabularx}
\end{center}

\section{Arquitetura}

\subsection{Visão geral}

A arquitetura é \textbf{serverless}, sem servidor próprio: toda a infraestrutura de backend é fornecida pela plataforma Firebase (Google Cloud). O sistema é uma aplicação web \textbf{SPA} (Single Page Application) que conversa diretamente com os serviços do Firebase pela internet.

\begin{center}
\begin{tikzpicture}[
  node distance=0.9cm and 2.2cm,
  box/.style={rectangle, rounded corners, draw, thick, minimum width=3.4cm, minimum height=0.9cm, align=center, fill=white},
  grupo/.style={rectangle, rounded corners, draw, thick, inner sep=14pt, fill=gray!5}
]
  \node[box] (spa) {Navegador (equipe)\\ \small React + TypeScript (SPA)};
  \node[box, below=of spa, xshift=-3.2cm] (auth) {Authentication\\ \small login e perfis};
  \node[box, below=of spa] (fs) {Cloud Firestore\\ \small banco NoSQL};
  \node[box, below=of spa, xshift=3.2cm] (st) {Firestore\\ \small arquivos em partes};
  \node[box, below=1.1cm of fs, xshift=3.2cm] (rules) {Security Rules\\ \small acesso por perfil};
  \node[grupo, fit=(auth)(fs)(st)(rules), label={[anchor=north]north:FIREBASE}] (fb) {};
  \draw[-{Latex}, thick] (spa) -- node[right, font=\small]{HTTPS / SDK} (fb.north -| spa);
\end{tikzpicture}
\end{center}

\subsection{Stack tecnológica}

\begin{center}
\begin{tabularx}{\textwidth}{l l X}
\toprule
\textbf{Camada} & \textbf{Tecnologia} & \textbf{Função} \\
\midrule
Front-end & React 19 + TypeScript + Vite & Interface web responsiva (computador e tablet) \\
Autenticação & Firebase Authentication & Login com e-mail/senha; perfis de acesso \\
Banco de dados & Cloud Firestore (NoSQL) & Alunos, turmas, oficinas, frequência, evolução \\
Arquivos & Cloud Firestore & Metadados e conteúdo em partes de até 600 KiB por documento \\
Segurança & Firestore Security Rules & Permissões por perfil (negar por padrão) \\
Ferramentas & Firebase CLI + Emulator Suite & Deploy, testes de regras e desenvolvimento local \\
Seed de dados & Node.js (script próprio) & População do banco com dados fictícios (LGPD) \\
\bottomrule
\end{tabularx}
\end{center}

\subsection{Decisões de arquitetura}

\textbf{Por que Firestore (NoSQL) e não um banco relacional?}
\begin{itemize}[nosep]
  \item Escala pequena (cerca de 5 profissionais e algumas dezenas de alunos) --- as limitações de consultas do NoSQL não impactam o sistema;
  \item Modelagem flexível: os requisitos ainda estão em definição com a Fundação, e documentos NoSQL absorvem mudanças de campos sem migrações;
  \item Autenticação, banco, arquivos e controle de acesso integrados em uma única plataforma, com plano gratuito;
  \item Regras de segurança declarativas (Security Rules), aplicadas no servidor do Firebase, independentes do front-end.
\end{itemize}

\textbf{Padrões adotados}
\begin{itemize}[nosep]
  \item \textbf{Negar por padrão (deny-by-default):} toda requisição é negada a menos que uma regra autorize explicitamente;
  \item \textbf{Dados sensíveis isolados:} a ficha médica fica em subcoleção própria, com regras mais restritivas (LGPD);
  \item \textbf{Dados que pertencem a um pai ficam embutidos ou em subcoleções} (ex.: presenças dentro da sessão de chamada), reduzindo o número de coleções e de consultas;
  \item \textbf{Datas puras como string \texttt{YYYY-MM-DD}} (evita problemas de fuso horário do Timestamp UTC) e data-hora como \texttt{Timestamp};
  \item \textbf{Auditoria padronizada:} documentos alteráveis possuem \texttt{criado\_em/criado\_por/atualizado\_em/atualizado\_por}.
\end{itemize}

\subsection{Segurança e LGPD}

\begin{itemize}[nosep]
  \item Perfis de acesso: \textbf{Direção}, \textbf{Pedagogia} e \textbf{Administrativo}, definidos no campo \texttt{papel} do usuário;
  \item Regras de leitura/escrita distintas por coleção (ex.: apenas a Direção apaga registros e gerencia usuários);
  \item Ficha médica: leitura restrita a Direção e Pedagogia; escrita apenas pela Direção;
  \item Nenhum dado pessoal real é usado em desenvolvimento --- apenas dados fictícios;
  \item Arquivos ficam no Firestore: metadados em \texttt{anexos/\{id\}} e conteúdo dividido em \texttt{chunks}, respeitando o limite de 1 MiB por documento. Não há dependência de Firebase Storage.
\end{itemize}

\section{Banco de Dados (Cloud Firestore)}

\subsection{Modelo de dados}

Firestore é um banco \textbf{NoSQL orientado a documentos}: os dados vivem em \textbf{coleções} (pastas) que contêm \textbf{documentos} (registros), identificados por IDs. Não há tabelas, joins nem schema fixo --- relacionamentos são representados por IDs e por subcoleções.

O modelo segue o princípio: \textbf{dados que só fazem sentido dentro de um pai não ganham coleção própria} --- viram subcoleção (listas independentes) ou campo embutido (poucos itens, limitados ao pai).

\subsection{Coleções e subcoleções}

\begin{verbatim}
usuarios/{uid}
alunos/{aluno_id}
  |-- ficha_medica/{ficha}        (subcoleção 1:1 -- dados sensíveis)
  |-- responsaveis/{responsavel_id}
  `-- documentos/{documento_id}
turmas/{turma_id}
oficinas/{oficina_id}
turma_oficinas/{turma_oficina_id}     (com horários embutidos)
matriculas/{matricula_id}
aluno_oficinas/{vinculo_id}
registros_atividade/{registro_id}
  `-- anexos/{anexo_id}
registros_evolucao/{registro_evolucao_id}
modelos_formulario/{modelo_id}
relatorios_pedagogicos/{relatorio_id}
  `-- anexos/{anexo_id}
sessoes_chamada/{sessao_id}           (com presenças embutidas)
contatos_familia/{contato_id}
relatorios_gerados/{relatorio_gerado_id}
configuracoes_sistema/geral
\end{verbatim}

\subsection{Principais atributos}

\subsubsection*{\texttt{usuarios/\{uid\}} --- ID igual ao UID do Firebase Authentication}

\begin{verbatim}
nome_completo: string          obrigatório
email: string                  obrigatório
papel: string                  obrigatório (direcao | pedagogia | administrativo)
ativo: boolean                 obrigatório
ultimo_login_em: Timestamp     opcional
criado_em / atualizado_em: Timestamp
\end{verbatim}

\subsubsection*{\texttt{alunos/\{aluno\_id\}}}

\begin{verbatim}
nome_completo: string          obrigatório
nome_busca: string             obrigatório (sem acentos, minúsculas -- auxilia busca)
data_nascimento: string        obrigatório (YYYY-MM-DD)
status: string                 obrigatório (ativo | inativo)
uso_imagem_autorizado: boolean obrigatório (LGPD)
consentimento_imagem_em: Timestamp (opcional)
caminho_foto: string           opcional (referência do arquivo no Firestore)
\end{verbatim}

\subsubsection*{\texttt{alunos/\{aluno\_id\}/ficha\_medica/\{ficha\}} --- subcoleção com regras restritivas}

\begin{verbatim}
tipos_deficiencia: array<string>      obrigatório
codigos_cid: array<string>            opcional
formato_ficha_medica: string          opcional (anexo | digital)
dados_ficha_medica: map               opcional
\end{verbatim}

\subsubsection*{\texttt{alunos/\{aluno\_id\}/responsaveis/\{responsavel\_id\}}}

\begin{verbatim}
nome_completo, parentesco, telefone,
telefone_alternativo (opcional), email (opcional), principal: boolean
\end{verbatim}

\subsubsection*{\texttt{alunos/\{aluno\_id\}/documentos/\{documento\_id\}}}

\begin{verbatim}
tipo_documento: string        (ficha_medica | ficha_matricula | autorizacao_imagem | outro)
armazenamento: "firestore", status: enviando | disponivel, total_partes
tamanho_bytes, nome_arquivo_original, tipo_conteudo, enviado_por, enviado_em
\end{verbatim}

\subsubsection*{\texttt{turmas/\{turma\_id\}}}

\begin{verbatim}
codigo, nome, periodo (manha | tarde | integral),
ids_usuarios_responsaveis: array<string>, data_inicio, data_fim (opcional), ativo
\end{verbatim}

\subsubsection*{\texttt{oficinas/\{oficina\_id\}}}

\begin{verbatim}
nome, categoria (musicalizacao, danca, culinaria, horta...),
descricao (opcional), ids_usuarios_responsaveis, ativo
\end{verbatim}

\subsubsection*{\texttt{turma\_oficinas/\{turma\_oficina\_id\}} --- vínculo N:N entre turma e oficina}

\begin{verbatim}
turma_id, oficina_id, ids_usuarios_responsaveis, ativo
horarios: array<map>          EMBUTIDO -- cada item:
  { dia_semana (1 a 7), hora_inicio (HH:mm), hora_fim, valido_de, valido_ate }
\end{verbatim}

\subsubsection*{\texttt{matriculas/\{matricula\_id\}} --- histórico de vínculo aluno-turma}

\begin{verbatim}
aluno_id, turma_id, data_inicio, data_fim (opcional)
status (ativo | concluido | cancelado)
\end{verbatim}

\subsubsection*{\texttt{aluno\_oficinas/\{vinculo\_id\}} --- plano individual de oficinas}

\begin{verbatim}
aluno_id, turma_oficina_id, data_inicio, data_fim (opcional)
status (ativo | concluido | cancelado)
nivel_suporte (opcional), consideracoes_pedagogicas (opcional), atribuido_por
\end{verbatim}

\subsubsection*{\texttt{registros\_atividade/\{registro\_id\}} --- registro diário por turma}

\begin{verbatim}
data_atividade: string        (YYYY-MM-DD)
turma_id, ids_turma_oficinas: array<string> (opcional)
usuario_responsavel_id, observacoes_gerais (opcional)
\end{verbatim}

\subsubsection*{\texttt{sessoes\_chamada/\{sessao\_id\}} --- chamada diária ou por oficina}

\begin{verbatim}
data_chamada: string          (YYYY-MM-DD)
turma_id, escopo (turma | oficina), turma_oficina_id (se oficina)
presencas: map                EMBUTIDO -- { aluno_id: presente | ausente | justificado }
observacoes: map              opcional -- { aluno_id: observação individual }
registrado_por
\end{verbatim}

\textit{As presenças ficam embutidas na sessão: um único documento por chamada permite montar a tela de chamada inteira com uma leitura, e torna impossível registrar a mesma presença duas vezes.}

\subsubsection*{\texttt{registros\_evolucao/\{registro\_evolucao\_id\}}}

\begin{verbatim}
aluno_id, data_registro (YYYY-MM-DD), profissional_usuario_id
area (cognitiva | motora | social | comunicacao | autonomia | outra)
observacao
\end{verbatim}

\subsubsection*{\texttt{modelos\_formulario/\{modelo\_id\}} --- formulários versionados}

\begin{verbatim}
nome, tipo_modelo (avaliacao_pedagogica), versao: number, ativo
campos: array<map> -- cada item: { chave, rotulo, tipo, obrigatorio, opcoes }
\end{verbatim}

\textit{Como o formulário oficial da Fundação ainda será definido, os modelos são versionados e as respostas ficam gravadas como map no relatório --- o sistema se adapta sem alteração de estrutura.}

\subsubsection*{\texttt{relatorios\_pedagogicos/\{relatorio\_id\}}}

\begin{verbatim}
aluno_id, tipo_relatorio (diario | semestral | anual)
periodo_inicio, periodo_fim (YYYY-MM-DD)
evolucao, dificuldades, avancos, encaminhamentos (opcionais)
usuario_responsavel_id, modelo_formulario_id + versao, respostas: map
status (rascunho | concluido)
\end{verbatim}

\subsubsection*{\texttt{contatos\_familia/\{contato\_id\}} --- contatos por faltas recorrentes}

\begin{verbatim}
aluno_id, responsavel_id (opcional)
datas_faltas_relacionadas: array<string> (YYYY-MM-DD)
contatado_em: Timestamp, meio_contato (telefone | whatsapp | email | presencial | outro)
observacoes, usuario_responsavel_id
\end{verbatim}

\subsubsection*{\texttt{relatorios\_gerados/\{relatorio\_gerado\_id\}} --- prestação de contas}

\begin{verbatim}
tipo_relatorio (acompanhamento_aluno | prestacao_contas)
periodo_inicio, periodo_fim, filtros: map (opcional)
status (processando | concluido | falhou), caminho_storage (PDF)
gerado_por, gerado_em
\end{verbatim}

\subsubsection*{\texttt{configuracoes\_sistema/geral}}

\begin{verbatim}
modo_chamada (diaria | oficina)
limite_faltas_consecutivas: number   (dispara alerta à equipe)
\end{verbatim}

\subsection{Relacionamentos}

\begin{center}
\begin{tabularx}{\textwidth}{l c X}
\toprule
\textbf{Origem} & \textbf{Relac.} & \textbf{Como é representado} \\
\midrule
Aluno & 1:1 & Ficha médica --- subcoleção \texttt{ficha\_medica} \\
Aluno & 1:N & Responsáveis --- subcoleção \texttt{responsaveis} \\
Aluno & 1:N & Documentos --- subcoleção \texttt{documentos} \\
Aluno & N:N & Turmas --- coleção \texttt{matriculas} (preserva histórico) \\
Turma & N:N & Oficinas --- coleção \texttt{turma\_oficinas} (horários embutidos) \\
Aluno & N:N & Oficinas --- coleção \texttt{aluno\_oficinas} \\
Turma & 1:N & Registros de atividade \\
Atividade / Relatório & 1:N & Anexos --- subcoleções \texttt{anexos} \\
Sessão de chamada & embutido & Presenças no map \texttt{presencas} da própria sessão \\
Aluno & 1:N & Contatos com a família \\
\bottomrule
\end{tabularx}
\end{center}

\subsection{Índices compostos declarados}

Definidos em \texttt{firestore.indexes.json} e aplicados automaticamente no deploy:

\begin{center}
\begin{tabularx}{\textwidth}{l X}
\toprule
\textbf{Coleção} & \textbf{Campos} \\
\midrule
alunos & status ASC, nome\_busca ASC \\
matriculas & aluno\_id ASC, status ASC \\
matriculas & turma\_id ASC, status ASC \\
turma\_oficinas & turma\_id ASC, ativo ASC \\
registros\_atividade & turma\_id ASC, data\_atividade DESC \\
registros\_evolucao & aluno\_id ASC, data\_registro DESC \\
relatorios\_pedagogicos & aluno\_id ASC, tipo\_relatorio ASC, periodo\_fim DESC \\
sessoes\_chamada & turma\_id ASC, data\_chamada DESC \\
sessoes\_chamada & turma\_oficina\_id ASC, data\_chamada DESC \\
contatos\_familia & aluno\_id ASC, contatado\_em DESC \\
aluno\_oficinas & aluno\_id ASC, status ASC \\
aluno\_oficinas & turma\_oficina\_id ASC, status ASC \\
modelos\_formulario & tipo\_modelo ASC, versao DESC \\
\bottomrule
\end{tabularx}
\end{center}

\subsection{Controle de acesso (Security Rules)}

Arquivo \texttt{firestore.rules} --- \textbf{nega tudo por padrão} e concede acesso por perfil:

\begin{center}
\small
\begin{tabularx}{\textwidth}{X l l}
\toprule
\textbf{Recurso} & \textbf{Leitura} & \textbf{Escrita} \\
\midrule
usuarios & autenticados & Direção \\
alunos & autenticados & Direção/Administrativo (criar/excluir); autenticados (atualizar) \\
ficha\_medica & Direção + Pedagogia & Direção \\
responsaveis / documentos & autenticados & Direção/Administrativo \\
turmas, oficinas, turma\_oficinas, matriculas & autenticados & Direção, Pedagogia ou Administrativo \\
aluno\_oficinas & autenticados & Direção/Pedagogia \\
registros\_atividade, evolução, anexos & autenticados & autenticados (criar/atualizar); exclusão só Direção \\
modelos\_formulario & autenticados & Direção/Pedagogia \\
relatorios\_pedagogicos & autenticados & Direção/Pedagogia \\
sessoes\_chamada (frequência) & autenticados & Direção, Pedagogia ou Administrativo \\
contatos\_familia & autenticados & Direção, Pedagogia ou Administrativo \\
relatorios\_gerados, configuracoes\_sistema & autenticados & Direção \\
\bottomrule
\end{tabularx}
\end{center}

\subsection{População e teste do banco (seed)}

O script \texttt{scripts/seed-firestore.mjs} popula o banco com \textbf{dados fictícios} (conforme LGPD, nenhum dado real), permitindo validar a modelagem e desenvolver as telas:

\begin{center}
\begin{tabularx}{\textwidth}{l X}
\toprule
\textbf{Comando} & \textbf{Função} \\
\midrule
\texttt{npm run seed} & Popula o banco (via acesso administrativo) \\
\texttt{npm run seed:reset} & Remove apenas os documentos criados pelo seed \\
\texttt{npm run seed:cleanup} & Limpeza completa, incluindo coleções legadas e subcoleções órfãs \\
\bottomrule
\end{tabularx}
\end{center}

O script valida o modelo ponta a ponta: alunos com ficha médica e responsáveis, turmas, oficinas com horários, matrículas, chamadas com presenças embutidas, contatos por falta, evolução, formulário versionado e relatório pedagógico.

\subsection{Pendências de validação com a Fundação}

\begin{itemize}[nosep]
  \item Campos definitivos do relatório pedagógico e formato oficial do formulário de avaliação;
  \item Formato da ficha médica (digitalizada, digitada ou ambos);
  \item Registro de frequência por turma/dia ou por oficina;
  \item Níveis de suporte/dependência dos alunos;
  \item Matriz de acesso aos dados médicos (participação de terapeutas);
  \item Quantidade de faltas consecutivas que dispara alerta;
  \item Estratégia de backup (o plano gratuito do Firebase não inclui backup diário automático).
\end{itemize}

\vspace{1cm}
\begin{center}
\small\textit{Documento gerado a partir da modelagem oficial do projeto (\texttt{modelagem.md}) e do código-fonte do repositório.}
\end{center}

\end{document}





